%% title: Adaptation is a classifier problem
%% target: Sections/Results.tex
%% anchor: Combined with the larger backbone and without oversampling (\cref{sec:results_longtail}), self-training gives B8 (0.647 and 0.452).
%% action: after
%% why: Adapting only the classifier on a frozen backbone captures most of self-training's gain in minutes; the staged recipe is better on the test fold and close under shift; balancing pseudo-labels harms a trainable backbone; a foreign backbone must be unfrozen. With the long-tail and hierarchy results, the interventions that matter act on the classifier.
%% source: paper/DRAFT.md §4.15, §4.14.3, §4.12 (centroids 0.29 vs 3.7 pt); journal/2026-08-06-adaptation-is-mostly-a-classifier-problem.md; journal/2026-08-27-the-noise-floor-does-not-transfer-across-training-regimes.md; data/rescore/probe-grouped.csv, probe-ho-grouped.csv (B1 0.4242/0.3561, B3rep5x 0.6309/0.4139, T2 0.6045/0.4338, T2b 0.6015/0.4191, F2 0.6019/0.3938, B9 0.5670/0.3730, G3 0.6181/0.4484, G3b 0.6195/0.4261, P4 0.4796/0.3568, P5 0.6223/0.4974)
%% check: Paste after step 20. | Numbers are recipe L before the rerun. | The staged in-distribution advantage (F2 0.9081 vs 0.9003) is unchanged in the current convention. | The centroid sentence could move to the appendix if space is short.
%% status: ready
\subsection{Adaptation is a classifier problem}\label{sec:results_classifier}
Training only the classifier of the frozen B1 backbone for 2 epochs on the same pseudo-labels (\cref{sec:selftraining}) raises the probe set from 0.424 to 0.605, 87\,\% of what retraining the whole network for 5 epochs achieves (0.631), and the held-out species from 0.356 to 0.434, above end-to-end self-training (0.414).
From a backbone trained without trap-like augmentation the result is almost the same (0.602), so augmentation and classifier adaptation are substitutes.
The full staged recipe, a representation trained without oversampling or augmentation followed by cRT and classifier adaptation (2 + 2 epochs on a frozen backbone), is 0.8 points better than end-to-end training on the test fold (0.9081 against 0.9003) and 2.9 points worse on the probe set (0.602 against 0.631), at a fraction of the cost.

Two limits bound the claim.
Balancing the pseudo-labels so that every species contributes equally harms a trainable backbone (end-to-end at 20\,M: $-6.4$ points on the probe set, $-4.1$ on held-out species), because balancing replicates the species with the fewest images, whose pseudo-labels are least reliable, and the backbone memorises them; on a frozen backbone its effect is comparable to the spread between identical runs (at 198\,M, $+2.4$ points on held-out species against a 2.2-point spread).
And a frozen backbone works only if it was trained with the head being refitted: BioCLIP-2's features, arranged for alignment with text, give 0.480 with frozen adaptation and 0.622 unfrozen (\cref{sec:results_contamination}).

With the hierarchy (coarse losses help, coarse parameters do not, \cref{tab:heads_recipeL}) and the long tail (rebalance the classifier, not the representation, \cref{sec:results_longtail}), this is the third question whose answer lies in the classifier.
For the margin head, the trained prototype matrix can even be replaced at inference by the mean training embedding of each class for 0.29 points (3.7 for the plain head).
The representation is the expensive and robust component; the classifier is the cheap one where interventions belong, and a new camera needs a new classifier, which takes minutes and no labels, rather than a new model.
